\section{Podsumowanie}
W ramach pracy udało się przeprowadzić badania z wykorzystaniem LIME, SHAP, PFI i Grad-CAM. Metody pozwalające na ich zastosowanie w przypadku danych tabelarycznych tłumaczyły wyniku uzyskane przez Random Forest, XGBoost'a i sieć neuronową, a obrazowe - VGG-16, ResNet50 i Xception. Przy CNN zastosowano \textit{fine tuning} zamrażając wagi bazowych modeli i dodając na ich końcu warstwy dostosowujące architekturę pod predykcję nowego problemu. W eksperymentach użyto 4 zbiorów danych, które zostały odpowiednio przetworzone. Proces ten obejmował uzupełnienie braków, usunięcie wartości odstających i duplikatów, normalizację (wymienione pozycje odnoszą się do danych tabelarycznych) oraz podział na podzbiory - treningowy i testowy. Zestawy pozwalały na przeprowadzenie zadań klasyfikacji - 3 binarnej i 1 wieloklasowej z czego dwa z nich składały się z rekordów, a dwa z obrazów.    

Zaprezentowane metody rzeczywiście pozwalają na wgląd w pracę modelu. Potrafią dostarczać wyjaśnienia bez wątpienia pomagające nam zrozumieć czym kieruje się badany system. Możemy za ich pomocą sprawdzić czy badana architektura osiągająca bardzo dobre rezultaty kieruje się poprawnymi, znanymi nam istotnymi czynnikami, zauważyć dzięki niej nowe zależności, których wcześniej nie dostrzegaliśmy, bądź nie braliśmy pod uwagę. A te modele, których wyniki są niezadowalające, zdiagnozować i zobaczyć dlaczego ich miary są tak mizerne.

LIME będący za każdym razem pierwszą omawianą metodą w przypadku danych tabelarycznych tworzył czytelny wykres wodospadowy. Patrząc na niego człowiek od razu wiedział, która z cech w jaki sposób wpływa na przewidywanie. Przedstawione na nim dane nie zlewają się ani nie zachodzą na siebie, dając przy tym większy wgląd w rozumienie badanego modelu. Wyniki dla każdej z próbek mogą zostać zebrane w celu stworzenia podsumowania w ich ogólną istotność na przestrzeni całego zbioru. Niestety tłumaczeniom ciężko jest zaufać ze względu na to, że mogą się one zmienić po ponownym uruchomieniu programu. Uzyskane wartości są jedynie relatywne, pozwalając tylko na interpretację porównawczą wzgledem siebie i nie pokazują rzeczywistego wpływu na predykcję, a wyłącznie w pewnym sensie przybliżony.
Wykorzystany do tłumaczeń decyzji klasyfikacyjnych zdjęć przedstawia konkretne, sprecyzowane obszary oznaczając je odpowiednim kolorem w zależności od ich wpływu na predykcję. Takie rozwiązanie jest bardzo precyzyjne, aczkolwiek metoda często wśród pierwszych najistotniejszych cech nie zwraca tej, która rzeczywiście powinna się tam znaleźć. Powoduje to konieczność jego odpowieniego dostrojenia, zwiększenia przestrzeni rozwiązań, bądź badanego sąsiedztwa.

SHAP pozwalał na zaprezentowanie wyników danych tabelarycznych w formie bardzo zmyślnie skonstruowanego wykresu wodospadowego. Początkowa problematyka związana ze zrozumieniem działania tak skomplikowanej metody wiązała się z poźniejszym przystępnym tłumaczeniem względem wartości średniej. Prezentowane wyniki dawały bezpośredni wgląd cech w ich wpływ na uzyskiwane prawdopodobieństwa lub \textit{log odds}. W jego ofercie znajdzie się wiele przydatnych prezentacji w postaci wykresów typu \textit{,,beeswarm''} czy \textit{,,force''}. Uzyskane przez niego tłumaczenia danych obrazowych pozostawiają jednak wiele do życzenia. Wskazane przez ową metodę obszary rozwiązań pokrywały zbyt duży obszar zdjęć szczycąc się niezwykle niską precyzją, dając wrażenie, prawie że wyboru losowego, gdyby nie wyraźne oznaczanie intensywnością barw miejsc szczególnie istotnych. Zaprezentowane wyjaśnienia bazują na kolorystyce, które poprzez skojarzenia tworzą nieskomplikowane w interpretacji wyniki. 

PFI pomimo bycia najprostszą metodą nie oferował tłumaczeń odbiegających jakościowo od jego konkurencji. Oprócz nieskomplikowanego działania jego wyniki były równie proste w interpretacji, dając wgląd na to, jak każda z cech wpływa bezpośrednio na dokładność modelu. Należy jednak pamiętać, że jest to metoda permutująca kolumny, przez co może ona tworzyć i poddawać ewaluacji dane nie posiadające odwzorowania w świecie rzeczywistym. Nie potrafi analizować danych będącymi obrazami, przez co możliwości jej wykorzystania są ograniczone, niemniej jednak tam gdzie można ją użyć radzi sobie dobrze. 

Grad-CAM, będący ostatnią badaną metodą w owej pracy, tworzył odznaczające się mapy cieplne. Nie znajdziemy na nich cech przeczących danej predykcji, a jedynie oznaczenia tych najistotniejszych pozytywnie miejsc. Uważam, że jest to pośredni wybór pomiędzy innymi omówionymi już metodami. Czasami tak jak LIME wskaże jeden konkretniejszy obszar, co prawda, o wiele mniej precyzyjny, a z drugiej strony zdarza mu się oznaczyć wiele fragmentów zdjęcia, jak ma to miejsce w przypadku SHAP, uznając ewidentnie zbyt wiele obszarów za ważnych.

Po przeanalizowaniu wad i zalet badanych metod, wydaje się, że pomimo problematyki związanej ze zrozumieniem działania najlepszym wyborem w przypadku danych tabelarycznych będzie SHAP. Jest on popularnym rozwiązaniem oferującym wiele możliwości wizualizacji otrzymanych wyników. Należałoby jednak wytłumaczyć grupie docelowej jego nieprostolinijne działanie. Dla przykładowych lekarzy, będacych bardzo zapracowanymi ludźmi, takie rozwiazanie mogłoby być nieodpowiednie. W przypadku tych danych zdaje się, iż najgorzej poradziła sobie metoda LIME. Głównie ze względu na nierelatywność jej miar lepsze prawdopodobnie byłoby wykorzystanie jej zamienników, w tym PFI.  

W kontekście zbiorów składających się z obrazów wygląda na to, że najlepiej prezentuje się LIME. Wyjątkowo niesatysfakcjonująco w obu zadaniach klasyfikacyjnych poradził sobie SHAP. W jednym z przypadków oznaczył praktycznie wszystkie obszary na zdjęciu jako istotne w predykcji. Grad-CAM wydaje się być od niego lepszą alternatywą, aczkolwiek, co widać szczególnie na zbiorze dotyczącym psów i kotów, zaznaczone przez niego obszary są zastanawiające.   

Wszystkie użyte metody są różne pod względem działania czy uzyskanych tłumaczeń, więc wybranie konkretnej z nich będzie zależeć od preferencji poszczególnych użytkowników oraz rodzaju zadania. W owej pracy starałam się przybliżyć ich działanie, a także mocne i słabe strony. Odpowiadając na wcześniej ustaloną tezę, jaką był wymóg sprawdzania systemów AI, to dalej się z nią zgadzam. Niektóre z tłumaczeń wydają się być nie do końca poprawne czy też trafione. Jeśli mielibyśmy powierzać zadania z krytycznych sektorów naszych żyć systemom używającym sztucznej inteligencji uważam, że musimy zapewnić ich jak największą niezawodność, którą osiągnąć możemy poprzez wykorzystanie metod XAI. 

Jedną z kluczowych pozycji, których zabrakło w tej pracy jest wykres ,,beeswarm'' LIME. Sama metoda koncentruje się na wyjaśnianiu lokalnym, czyli pojedynczych próbek, przez co biblioteka ją implementująca nie zawiera już przygotowanych sposobów na uzyskanie jej wyników dla całej populacji zbioru. Problem ten można rozwiązać agregując samemu dane dla każdej z próbek, co zostało zrobione, jednak wśród wyników zabrakło wykresu typu ,,rój pszczół''. Wydaje się, że metody LIME i SHAP są do siebie bardzo podobne. Czyni je to konkurencyjnymi rozwiązaniami. Inna sytuacja jest w przypadku PFI, podejście zdaje się nie wykluczające z innymi. Prawdopodobnie możnabyłoby stworzyć jedno rozwiązanie łączące LIME lub SHAP z PFI. Innym problemem jest nienajprzejrzystrze wizualizacje wykresu ,,beeswarm'' SHAP dla problemu klasyfikacji wieloklasowej. Potrzebuje on wektora składającego się z jednego wymiaru - zapisuje się tylko wskazane przez model klasy predykcyjne. Powoduje to interpretację w stylu: ,,Analizując na wykresie zmienną odpowiadającą wiekowi pacjenta zauważamy znajdują się dwa skupiska niebieskich kropek - jedno z lewej strony od wartości średniej dla predykcji drugie zacznie z jej prawej. Oznacza to, że niski wiek respondenta dla jakiejś klasy zmniejsza prawdopodobieństwo danej predykcji, a dla innej w zauważalym stopniu je zwiększa.'' Brakuje tu informacji, których klas dotyczy analiza. Najpewniej możliwe jest rozszerzenie wizualizacji o ową treść, np. poprzez stworzenie tła za punktami. Owe propozycje dotyczyły części tabelarycznej. W przypadku zbiorów z obrazami wyniki Grad-CAM zdawały się być nie do końca sprecyzowane. Jest to metoda, która działa w sposób bardzo zintegrowany z analizowanym modelem, czyniąc ją w ten sposób unikatową względem innych użytych w tej pracy rozwiązań. Dobrze byłoby gdyby powstały jakieś sposoby jej dostrojenia i zwiększenia dokładności.  